\chapter{市场重置：8月10日收盘，而非8月11日盘前}

\section{四只原主标的位置}

市场数据并不产生目标价，但会改变旧标签是否还能成立。四只原主标的中，只有恒逸已从低位直接切换到高位盈利验证；其余三只仍处一年区间低段，却不能因为回撤而跳过利润和现金门槛。

\begin{exhibitbox}[图表3｜价格位置与短期变动]
\small
\begin{tabularx}{\textwidth}{L{1.1cm}L{1.55cm}R{1.0cm}R{1.25cm}R{1.35cm}R{1.2cm}X}
\toprule
\textbf{代码} & \textbf{标的} & \textbf{收盘} & \textbf{一年位置} & \textbf{距一年高点} & \textbf{20日涨跌} & \textbf{含义} \\
\midrule
000703 & 恒逸石化 & 18.18 & 99.96\% & -0.03\% & +37.52\% & 价格已先行；不能沿用低位标签。 \\
002379 & 宏桥控股 & 19.76 & 30.94\% & -38.79\% & +24.28\% & 仍属低区间，但已出现反弹。 \\
300014 & 亿纬锂能 & 55.22 & 22.45\% & -39.16\% & +6.73\% & 低区间只是等待经营验证的背景。 \\
002497 & 雅化集团 & 19.09 & 32.67\% & -38.46\% & +5.35\% & 周期价格仍在验证阶段。 \\
\bottomrule
\end{tabularx}
\normalsize
\sourcenote{腾讯复权收盘用于一年位置；搜狐日线用于成交额。}
\end{exhibitbox}

图表3的决策含义是：恒逸的价格门槛已经失效，宏桥、亿纬和雅化的价格位置只允许其进入下一份披露的验证序列。没有一项价格位置指标被用作支持位、阻力位或估值目标。

\section{成交与流动性边界}

8月10日成交额分别为恒逸CNY20.74亿、宏桥CNY12.51亿、亿纬CNY21.94亿、雅化CNY9.46亿。报告不收集席位、北向、融资余额或热钱标签，因此不对机构资金态度作因果判断。价格区间只用于限制叙事，不能代替估值或交易支持阻力位。

\section{历史高空间样本的市场含义}

九安医疗位于一年区间65.29\%，不再是干净的低位样本；江波龙位于51.26\%，5日上涨28.99\%，高波动和较大成交额更需要现金验证。恩捷的CNY103外部参照从CNY55.95计算仅约84\%空间，已不构成翻倍。
